% Generated by tools/edn_latex.py from reports/edn.md.
% Do not edit: change reports/edn.md and run `make edn`.
\ednentry{
  id = {EDN-78},
  title = {The hyperparameter search extends EDN-73's protocol to both LightGBM variants and \texorpdfstring{\texttt{min\_\allowbreak{}child\_\allowbreak{}samples}}{min\_child\_samples}, and measures and tracks every fit},
  date = {2026-10-06},
  milestone = {M3: Quality Assurance},
  topic = {Hyperparameter Tuning, Experiment Tracking, Energy Efficiency, Requirements},
  participants = {@kadameit, who decided, with AI assistance},
  decision = {The LightGBM hyperparameters are searched by \texttt{recommenditos/\allowbreak{}modeling/\allowbreak{}tune.\allowbreak{}py}, outside \texttt{dvc repro}, under the protocol EDN-73 wrote down, extended in four ways.
Both variants are swept, \texttt{lgbm-\allowbreak{}basic} and \texttt{lgbm-\allowbreak{}extended}, on EDN-73's grid: \texttt{learning\_\allowbreak{}rate} 0.025, 0.05 and 0.1 against \texttt{num\_\allowbreak{}leaves} 31, 63, 127 and 255, every point at a ceiling of 20,000 trees so early stopping decides it.
A second grid follows for each variant, over \texttt{min\_\allowbreak{}child\_\allowbreak{}samples} 5, 10, 20, 50 and 100, with \texttt{learning\_\allowbreak{}rate} and \texttt{num\_\allowbreak{}leaves} fixed at the first grid's winner, which is the committed point when nothing was adopted; \texttt{min\_\allowbreak{}child\_\allowbreak{}samples} becomes an optional LightGBM key at LightGBM's own default of 20, the value every committed model already used, so \texttt{params.\allowbreak{}yaml} names it only if a search adopts another value.
The adoption rule is EDN-73's, unchanged: a point replaces the reference only when the simultaneous 95\,\% interval of its validation L1 difference, from the paired bootstrap resampled by seller group and widened by the max-statistic bootstrap over all the grid's points, lies entirely below zero. The search never reads the test or calibration split.
Every fit is measured with CodeCarbon and logged as its own MLflow run in \texttt{recommenditos-\allowbreak{}price-\allowbreak{}tuning}, tagged with its sweep, \texttt{sweep\_\allowbreak{}role=\allowbreak{}point} and its variant, carrying its validation metrics, fit time and energy; one \texttt{sweep\_\allowbreak{}role=\allowbreak{}summary} run per sweep carries the totals and the verdict.
A sweep's record -- its results as JSON, the same as a readable table, and CodeCarbon's row for every fit as a CSV -- is kept in MLflow only, as artifacts under \texttt{results/\allowbreak{}} of the sweep's summary run, \texttt{sweep.\allowbreak{}json}, \texttt{sweep.\allowbreak{}txt} and \texttt{emissions.\allowbreak{}csv} for a sweep run from now on; nothing of a search is committed.
The four sweeps of \#88 ran before this, so their summary runs carry the JSON and the table under the names the module then used, \texttt{results/\allowbreak{}<sweep id>.\allowbreak{}json} and \texttt{.\allowbreak{}txt}, and their CodeCarbon CSV, until then committed, was uploaded to the same run as \texttt{results/\allowbreak{}emissions.\allowbreak{}csv}.
The search's energy is reported as its own total there, never in \texttt{reports/\allowbreak{}emissions/\allowbreak{}<variant>.\allowbreak{}csv}, which belongs to the fit of the final configuration.
NFR-10 is reworded to match: the 15-minute budget applies to fitting the final chosen configuration, and a search runs separately, with its energy recorded per fit and reported apart from the final fit's.},
  rationale = {\emph{Alternatives considered:}
\begin{itemize}[leftmargin=*, topsep=2pt]
\item \textbf{Option A (chosen): extend EDN-73's protocol, and build its decision step into the package.}
\newline Pros: the first grid for \texttt{lgbm-\allowbreak{}basic} repeats EDN-73's twelve fits exactly, so the new tool is checked against recorded evidence before anything new is read from it; the simultaneous interval still guards against the winner's curse within each grid; the cost is bounded and small, about 34 fits in all, and EDN-73's twelve \texttt{lgbm-\allowbreak{}basic} fits took 127 s of fitting; and issue \#88's acceptance criteria -- a tagged run per configuration with its metric, fit time and energy, the search's emissions apart, NFR-10 reworded -- are met by code that runs every time rather than by hand once.
\newline Cons: on \texttt{lgbm-\allowbreak{}basic} the likely outcome is EDN-73's again, ``keep'', so the model card may still have to say the model is untuned; one validation split stays the only source of selection noise; the two grids are run one after the other, so an interaction between \texttt{min\_\allowbreak{}child\_\allowbreak{}samples} and the first grid's knobs is not searched; nothing corrects across the two grids or the two variants, because the max-statistic interval is taken per sweep, and the second grid is chosen after the first on the same validation split; and \texttt{feature\_\allowbreak{}fraction} and L2 regularisation are still not parameters.
\item \textbf{Option B: the issue as written -- grid search, and the best validation point goes into \texttt{params.\allowbreak{}yaml}.}
\newline Pros: the least machinery, it matches the issue text, and the model can be called tuned.
\newline Cons: it adopts noise. EDN-73 measured exactly this case: the lowest validation L1 of its grid, \texttt{(0.\allowbreak{}025, 63)}, is indistinguishable from the committed point even point by point, and the two metrics disagree on which point is best. It would also overturn the adoption rule EDN-73 decided, move every number in the report for no measurable gain, and with the issue's 0.03 the grid would no longer be comparable with EDN-73's measured 0.025.
\item \textbf{Option C: random or Optuna search over more knobs, selected by seller-grouped k-fold cross-validation on train plus validation.}
\newline Pros: cross-validation is a less noisy selection than one split, and a larger space might hold a real gain.
\newline Cons: a new dependency and new code; about five times the fits per point, so five times the energy, for a region EDN-73 found flat; early stopping then needs a validation fold inside every fold; and it departs from the problem specification's ``the validation set is used for early stopping and tuning''. EDN-73 left it for the day a criterion depends on a fraction of a point, and none does.
\item \textbf{Option D: close \#88 as covered by EDN-73.}
\newline Pros: no work and no compute.
\newline Cons: \texttt{lgbm-\allowbreak{}extended} would never have been swept and \texttt{min\_\allowbreak{}child\_\allowbreak{}samples} never examined, although many \texttt{model} values occur only once; EDN-73's screening runs carry no energy figure per point; and NFR-10 would still say ``no hyperparameter search'' with no rule for where a search's energy goes. None of \#88's acceptance criteria would be met.
\item \textbf{Where a sweep's record is stored (decided 2026-10-09, after the search had run).} The pipeline's own \texttt{reports/\allowbreak{}emissions/\allowbreak{}<variant>.\allowbreak{}csv} stays git-tracked whichever is chosen, because the \texttt{compare-\allowbreak{}energy} stage reads it; this is about the search's record only.
\par
\begin{itemize}[leftmargin=*, topsep=2pt]
\item \textbf{Option A: the results and CodeCarbon's CSV committed under \texttt{reports/\allowbreak{}tuning/\allowbreak{}}, as well as the runs in MLflow.}
\newline Pros: the numbers can be checked from the repository without access to the DagsHub tracking server, and it follows the convention of \texttt{reports/\allowbreak{}analysis/\allowbreak{}}, whose results files are committed beside their scripts.
\newline Cons: the same data in two places, which can disagree; and a commit is needed between sweeps, because a sweep's uncommitted files make every later run \texttt{git\_\allowbreak{}dirty} (EDN-74).
\item \textbf{Option B (chosen): MLflow only, with CodeCarbon's CSV uploaded as an artifact of the summary run.}
\newline Pros: one source of truth; the closest to the course demo, which logs CodeCarbon's row to MLflow as metrics and parameters and commits no emissions file; no commit between sweeps; and the raw CSV keeps how each figure was produced -- CodeCarbon's version, the CPU and RAM power columns, the operating system -- which matters because EDN-69 calls the figure an estimate, and an estimate is checked by its method.
\newline Cons: the numbers this entry and the model card quote can be checked only with access to the DagsHub tracking server.
\item \textbf{Option C: B without the upload, exactly as the course demo does, with the CSV discarded.}
\newline Pros: the least to store, and nothing beyond what the demo does.
\newline Cons: rejected, because it loses the columns that say how each figure was produced, which the per-run metrics and parameters carry only in part.
\end{itemize}
\end{itemize}
\par
\emph{Why:} EDN-73 already answered how a value may be adopted; what \#88 added was scope and record-keeping, so the protocol was extended rather than replaced.
Two parts of the issue were stale against EDN-73, which landed after it was written: it proposed \texttt{learning\_\allowbreak{}rate} 0.03 where EDN-73 measured 0.025, and it said the winning values go into \texttt{params.\allowbreak{}yaml}, where EDN-73 adopts only a value whose simultaneous interval lies below zero.
0.025 was kept because a repeat of EDN-73's grid is the one check of the new code against numbers already on record, and because 0.03 lies inside the 0.025 to 0.05 range EDN-73 found flat.
\texttt{min\_\allowbreak{}child\_\allowbreak{}samples} is the knob the issue named as optional, and the data gives it a reason: a leaf of fewer rows can fit one listing of a \texttt{model} seen once.
Making it optional at LightGBM's default rather than writing \texttt{min\_\allowbreak{}child\_\allowbreak{}samples: 20} into \texttt{params.\allowbreak{}yaml} keeps every committed parameter value as it is, and a test holds that leaving it out and setting it to 20 fit the same booster.
The second grid holds the first grid's winner fixed because a joint grid would be 60 fits per variant instead of 5, for an interaction nothing so far suggests.
\par
The decision step became a module of the package, run from the command line, rather than staying an analysis script, because \#88 asks for every searched configuration to be a tracked run with its energy, and that is what \texttt{train} already does for one fit: the search reuses its \texttt{read\_\allowbreak{}matrices}, \texttt{fit\_\allowbreak{}variant}, \texttt{energy.\allowbreak{}measure} and run logging, so a search run and a pipeline run are measured and named alike.
It is not a DVC stage, because a search's result is a verdict a person acts on through a reviewed change to \texttt{params.\allowbreak{}yaml}, not an artefact a later stage consumes, and its cost scales with the grid rather than with the product.
For the same reason NFR-10's 15 minutes apply to the final fit: the budget exists to keep retraining the product cheap, and a search is run once per decision, not once per retraining.
The search's energy is not written beside the variants' CSVs, because each of those is the record of the fit that produced its model, and because \texttt{compare-\allowbreak{}energy} declares the whole \texttt{reports/\allowbreak{}emissions} directory as a dependency, so a search writing there would make that stage stale.
Where it is kept instead is the sub-decision above: MLflow alone, as artifacts of each sweep's summary run, so a sweep refuses to start when tracking is not configured, opens its runs so that a server it cannot reach fails it rather than silently dropping them, and refuses a sweep id the experiment already holds runs for.
Queued \texttt{dvc exp} screening stays available as EDN-73 describes, and is now optional, because the decision step measures and logs every fit itself.
EDN-73's bootstrap moved into the module unchanged: \texttt{tuning\_\allowbreak{}sweep.\allowbreak{}py} now imports it and computes bit-identical draws and multipliers, so EDN-73's recorded evidence still reproduces.
\par
\textbf{Results (2026-10-06).} All four sweeps kept their reference, \texttt{learning\_\allowbreak{}rate} 0.05, \texttt{num\_\allowbreak{}leaves} 63 and \texttt{min\_\allowbreak{}child\_\allowbreak{}samples} 20, so nothing is adopted: \texttt{params.\allowbreak{}yaml} and \texttt{dvc.\allowbreak{}lock} are unchanged, and the committed test metrics and gate verdict, computed once by the committed pipeline run, are those of the configuration the search kept; the search never read the test split.
The search ran from clean commits on a 13th Gen Intel Core i7-1360P laptop (Windows, CodeCarbon 3.3.1's estimate, \texttt{num\_\allowbreak{}threads} 1), into the DagsHub MLflow experiment \texttt{recommenditos-\allowbreak{}price-\allowbreak{}tuning}; every one of the 34 fits was ended by early stopping, not by the 20,000-tree ceiling.
\par
\begin{itemize}[leftmargin=*, topsep=2pt]
\item \texttt{88-\allowbreak{}lgbm-\allowbreak{}basic-\allowbreak{}lr-\allowbreak{}leaves} reproduces EDN-73's twelve fits exactly: the same best round, validation L1 and validation MdAPE at every point, and the same per-point and simultaneous intervals, so the new module is checked against the recorded evidence, and only the fit times differ, because the machine does.
Its lowest validation L1 is again \texttt{(0.\allowbreak{}025, 63)}, 0.00031 below the reference with a simultaneous interval of -0.00109 to +0.00047, and only \texttt{(0.\allowbreak{}1, 127)} and \texttt{(0.\allowbreak{}1, 255)} are measurably worse; validation MdAPE spans 6.88\,\% to 7.04\,\%.
\item \texttt{88-\allowbreak{}lgbm-\allowbreak{}extended-\allowbreak{}lr-\allowbreak{}leaves}: the lowest validation L1 is \texttt{(0.\allowbreak{}025, 127)}, 0.00040 below the reference with a simultaneous interval of -0.00150 to +0.00070.
\texttt{(0.\allowbreak{}025, 255)} is the case the rule exists for: it is 0.20 pp better in validation MdAPE, 6.25\,\% against 6.45\,\%, with a per-point interval of -0.385\,\% to -0.060\,\%, entirely below zero, and a search that took the best validation MdAPE would have adopted it.
But MdAPE is not the selection metric, its L1 difference, -0.00039, has a per-point interval of -0.00146 to +0.00069 that the split cannot resolve, and even its MdAPE difference, widened to hold for all eleven points at once, spans zero at -0.427\,\% to +0.028\,\%: the best of eleven noisy comparisons looks better than it is, which is the winner's curse.
Five points are measurably worse under the simultaneous L1 interval: all four at \texttt{learning\_\allowbreak{}rate} 0.1, whose simultaneous intervals start at +0.00037 or higher, and \texttt{(0.\allowbreak{}05, 31)}, at +0.00003 to +0.00243; validation MdAPE spans 6.25\,\% to 6.69\,\%.
\item \texttt{88-\allowbreak{}lgbm-\allowbreak{}basic-\allowbreak{}min-\allowbreak{}child-\allowbreak{}samples}: the lowest validation L1 is \texttt{min\_\allowbreak{}child\_\allowbreak{}samples} 10, 0.00043 below the reference with a simultaneous interval of -0.00124 to +0.00038, and 100 is measurably worse, at +0.00010 to +0.00210; validation MdAPE spans 6.92\,\% to 7.03\,\%.
\item \texttt{88-\allowbreak{}lgbm-\allowbreak{}extended-\allowbreak{}min-\allowbreak{}child-\allowbreak{}samples}: the reference, 20, has the lowest validation L1 of the five, no point is resolved either way under the simultaneous interval, and validation MdAPE spans only 6.45\,\% to 6.49\,\%.
\end{itemize}
\par
The search cost 34 fits, 391.3 s of fitting (1,069.0 CPU-s), 3.6071 Wh and 627.81 mg CO2eq: 17 fits, 112.3 s, 1.0864 Wh and 189.09 mg for \texttt{lgbm-\allowbreak{}basic}, and 17 fits, 279.1 s, 2.5207 Wh and 438.73 mg for \texttt{lgbm-\allowbreak{}extended}, each sweep's totals on the ``search totals'' line of its table.
On the search's own machine that is about 20 times the energy of one fit of \texttt{lgbm-\allowbreak{}basic}'s committed configuration (0.0549 Wh) and 16 times one of \texttt{lgbm-\allowbreak{}extended}'s (0.1584 Wh); its fit times are not comparable with the ladder's energy table in the model card, which was measured on another laptop, an i5-10210U.
Because nothing was adopted, the final configuration is the committed one, and its fit time against NFR-10's 15 minutes is the one already on record: 8.3 s for \texttt{lgbm-\allowbreak{}basic} and 27.9 s for \texttt{lgbm-\allowbreak{}extended} in the tracked run of 2026-10-06; the search itself sits outside that budget, as the reworded NFR-10 says.
As option A's cons anticipated, the verdict on \texttt{lgbm-\allowbreak{}basic} is EDN-73's again, and the model card now says the values were searched and kept rather than that they are untuned.},
  ai = {Information seeking, Alternative generation, Alternative assessment, Recommendation},
  response = {Accepted with modifications},
  assessment = {AI compared issue \#88 with EDN-73 and the code, and pointed out where the issue had become stale: it proposed 0.03 where EDN-73 measured 0.025, and it said the winning values go into \texttt{params.\allowbreak{}yaml} where EDN-73 adopts only a value its rule resolves.
It laid out options A to D with their pros and cons and recommended option A, a second grid over \texttt{min\_\allowbreak{}child\_\allowbreak{}samples} with the first grid's winner fixed, recording the method decision as a new entry extending EDN-73, and \texttt{learning\_\allowbreak{}rate} 0.025.
Kevin weighed the options and accepted all four recommendations before the search was implemented or run.
The NFR-10 rewording and reporting the search's energy apart were not separate choices: they follow from the scope of issue \#88.
AI also pointed out that the search's emissions could not go under \texttt{reports/\allowbreak{}emissions/\allowbreak{}} without making the \texttt{compare-\allowbreak{}energy} stage stale.
On where a sweep's record should live, AI first recommended option A of that sub-decision, committing it under \texttt{reports/\allowbreak{}tuning/\allowbreak{}} beside MLflow, so the numbers could be checked without an account; on Kevin's question it checked the course demo repository and found that the demo keeps its emissions in MLflow only.
Kevin chose MLflow only, with CodeCarbon's CSV uploaded as an artifact, which is the modification.},
  aievidence = {Claude Code session of 2026-10-06 on issue \#88: the comparison with EDN-73 and options A to D with a recommendation, put to Kevin before the search was implemented; Claude Code session of 2026-10-09: the options for where a sweep's record is stored, with a recommendation for option A, the check of the course demo repository Kevin asked for, and his choice of option B.},
  evidence = {\href{https://github.com/mlops-2627q1-mds-upc/MLOPS_Recommenditos/issues/88}{issue \#88}; \href{https://github.com/mlops-2627q1-mds-upc/MLOPS_Recommenditos/blob/main/recommenditos/modeling/tune.py}{\texttt{recommenditos/\allowbreak{}modeling/\allowbreak{}tune.\allowbreak{}py}} and \href{https://github.com/mlops-2627q1-mds-upc/MLOPS_Recommenditos/blob/main/tests/test_tune.py}{\texttt{tests/\allowbreak{}test\_\allowbreak{}tune.\allowbreak{}py}}; \href{https://github.com/mlops-2627q1-mds-upc/MLOPS_Recommenditos/blob/main/reports/analysis/tuning_sweep.py}{\texttt{reports/\allowbreak{}analysis/\allowbreak{}tuning\_\allowbreak{}sweep.\allowbreak{}py}}, now importing the module's statistics; the search's results in the \href{https://dagshub.com/recommenditos/MLOPS_Recommenditos.mlflow}{DagsHub MLflow server}, experiment \texttt{recommenditos-\allowbreak{}price-\allowbreak{}tuning}: the sweeps \texttt{88-\allowbreak{}lgbm-\allowbreak{}basic-\allowbreak{}lr-\allowbreak{}leaves}, \texttt{88-\allowbreak{}lgbm-\allowbreak{}extended-\allowbreak{}lr-\allowbreak{}leaves}, \texttt{88-\allowbreak{}lgbm-\allowbreak{}basic-\allowbreak{}min-\allowbreak{}child-\allowbreak{}samples} and \texttt{88-\allowbreak{}lgbm-\allowbreak{}extended-\allowbreak{}min-\allowbreak{}child-\allowbreak{}samples}, each found with the filter \texttt{tags.\allowbreak{}sweep =\allowbreak{} \textquotedbl{}<id>\textquotedbl{}} and its record under \texttt{results/\allowbreak{}} of the run with \texttt{tags.\allowbreak{}sweep\_\allowbreak{}role =\allowbreak{} \textquotedbl{}summary\textquotedbl{}}, compared with EDN-73's \href{https://github.com/mlops-2627q1-mds-upc/MLOPS_Recommenditos/blob/main/reports/analysis/tuning_sweep_results.txt}{\texttt{tuning\_\allowbreak{}sweep\_\allowbreak{}results.\allowbreak{}txt}}; \href{https://github.com/mlops-2627q1-mds-upc/MLOPS_Recommenditos/blob/main/docs/docs/pipeline.md}{pipeline docs}, ``Tuning a hyperparameter''; \href{https://github.com/mlops-2627q1-mds-upc/MLOPS_Recommenditos/blob/main/docs/docs/specification.md}{specification} and \href{https://github.com/mlops-2627q1-mds-upc/MLOPS_Recommenditos/blob/main/docs/docs/requirements.md}{requirements} NFR-10; \hyperref[EDN-73]{EDN-73}, \hyperref[EDN-69]{EDN-69}, \hyperref[EDN-70]{EDN-70}.},
}
